\scanpage{101}\mspage{97}
Dans ce cas les correspondances \unsure{de|des} \unsure{disjoncteurs|disponcteurs} sont les corr. ensemblistes, elles forment l'ensemble $\unsure{\mathcal P|\mathscr P|P}(E\times E')$, qui est lui-même un \deleted{disjoncteur} \unsure{spacios|spatios}, associé au disjoncteur trivial $E\times E'$. Il serait \unsure{tentant} alors, dans le cas général, de vouloir interpréter les correspondances de disjoncteurs \inserted{de $\underline E$ à $\underline{E'}$} comme étant les éléments du spatios associé au ``disjoncteur produit'' $E\times E'$, avec
\[
 (x,x')\mathrel{|\,\unsure{\supset|\circ|o|0}\,|}(y,y')
 \overset{\text{\unsure{def|déf}}}{\Longleftrightarrow}
 x\mathrel{|\,\circ\,|}x'\quad\text{\uline{ou}}\quad y\mathrel{|\,\circ\,|}y'.
\]

On vérifie en effet que les \deleted{supports} $R\in\Sigma_{E\times E'}$ sont dans l'ens. noté $\operatorname{Corr}^{\unsure{\times|\star|\kappa}}(\underline E,E')$ (p. 94) i.e. $\deleted{\unsure{R|\mathcal R}}\forall \deleted{\unsure{\alpha|x}}x\in E,\ R(x)\in\Sigma_{E'}^{\unsure{7|\urcorner|\prime|?}}$, et $\forall x'\in E',\ {}^{t}R(x')\in\Sigma_E$. Mais il ne semble pas que \deleted{\unsure{tout|tait}} l'inverse soit vrai, \unsure{\uline{sauf}|sauf} dans le cas où $\underline E$ \uline{ou} $\underline{E'}$ est \deleted{\unsure{le|l'}}un disjoncteur trivial.

Mais si ce n'était pas vrai en général, la question demeure si l'ens. des \deleted{\unsure{disjo|dispo}} corr. de $\underline E$ à $\underline{E'}$, ou \inserted{encore} des \deleted{\unsure{morphismes}}\inserted{\unsure{corr..|corr.}} de spatios $\Sigma\unsure{\mbox{---}|\longrightarrow}\Sigma'$, forme un spatios. La relation d'ordre suggérée par le cas \unsure{d'ensembliste|ensembliste} (disj. triviaux) est simplement
% archive sheet number: 100
